\documentclass[11pt]{article}
\usepackage[utf8]{inputenc}
\usepackage[T1]{fontenc}
\usepackage{lmodern}
\usepackage[margin=1in]{geometry}
\usepackage{amsmath,amssymb}
\usepackage{textcomp}
\usepackage[dvipsnames]{xcolor}
\usepackage{enumitem}
\usepackage{parskip}
\usepackage{titlesec}
\usepackage[colorlinks=true,linkcolor=blue,citecolor=blue,urlcolor=blue]{hyperref}
\usepackage[round]{natbib}

\titleformat{\section}{\large\bfseries}{\thesection}{1em}{}
\titleformat{\subsection}{\normalsize\bfseries}{\thesubsection}{1em}{}
\titleformat{\subsubsection}{\normalsize\bfseries\itshape}{}{0em}{}

\title{\textbf{Manufacturing Ground Truth:\\ Mechanically Verifiable Failure Attribution via Controlled Component Substitution}}
\author{}
\date{}

\newcommand{\ph}{\textit{[placeholder --- fill in after experiments are run]}}

\begin{document}
\maketitle
\vspace{-2em}

\begin{quote}
\small\textit{Assembled from the six confirmed stage outputs of an \texttt{auto-research-discussion} run (seed topic B2.1), post research-taste audit (\textcolor{ForestGreen}{\textbf{PASS}}, no reframe needed). Revised 2026-08-15 per \texttt{review\_2026-08.md}: repositioned against Who\&When Pro (arXiv 2607.09996; concede-and-claim, certification-validity headline), bypass control arm added, label-survival statistic named, walking-skeleton pilot pre-registered. Status: \textbf{specification / paper draft --- no experiments have been run.} All numbers below are pre-registered design targets, not results.}
\end{quote}

\begin{abstract}
Failure-attribution methods for multi-agent and tool-use LLM systems are proliferating, and their evaluations fall into two families: agreement with human-annotated or single-culprit blame labels, and --- most recently --- injection-built corpora that replay a successful prefix, insert a fault, and declare the injected step a ``golden label'' by convention. Neither family verifies the causal premise it rests on: a method can produce a plausible blame assignment that is symptom-matching rather than causally correct, and an injected fault can be followed by a failure the injection did not actually cause. Whether ``failure follows the substituted component'' is even a reliable premise has never itself been measured: intervention-based auto-debugging, counterfactual label generation, and convention-labeled injection corpora all lean on it without testing its reliability across a plan, tool, and observation taxonomy. We build a controlled-substitution harness that checkpoints a reproducibly-succeeding agent run, injects one fault into one component while holding the other two fixed, and certifies a ground-truth label only when the outcome flips against a matched control lattice --- a class-matched no-op substitution, an identity replay, and a fault-bypass arm --- ruling out the substitution act, ordinary stochasticity, or the instrumentation itself as false-positive sources. The primary result is a validity measurement no existing corpus reports: per fault class, whether certified ground truth is obtainable at all, together with the fraction of convention-accepted labels the certification gate rejects. Secondarily, every candidate attribution method is applied uniformly to the certified corpus and scored against this ground truth, with a budget-matched control separating causal-intervention access from raw compute, measuring how far current methods fall short of the certified ceiling.
\end{abstract}

\section{Introduction}

Every current method for assigning blame to a specific component of a failing tool-use agent pipeline validates that blame against a human annotator's judgment call, an assumption that exactly one step is decisively responsible, or --- in the newest injection-built corpora --- the convention that an injected fault is the decisive error simply because it was injected. None of those checks is a certified causal ground truth, so a plausible-looking blame assignment may be symptom-matching rather than causally correct, and a convention-assigned label may name a fault the failure never actually tracked. Tool-use LLM agents increasingly carry out long, multi-step tasks that interleave a plan, an executor that acts on that plan, and a stream of tool calls and observations that feed back into what happens next; when such a pipeline fails, an operator needs to know which stage to fix --- a bad plan, a broken tool integration, or a misread observation call for different repairs --- and a growing line of work has proposed automated methods to answer exactly that question. A method can match what a human would have said, or correctly name the one step a human labeled decisive, while still being wrong about which component actually \emph{caused} the failure --- the two only coincide if human judgment and single-culprit structure are themselves reliable causal readings, and nothing in the current evaluation landscape checks whether they are. This matters beyond bookkeeping: intervention-based auto-debugging tools and counterfactual label-generation pipelines already build on a version of the premise that ``failure tracks the substituted component,'' rewinding to a suspected step, editing it, and replaying to see if the outcome changes. If that premise holds only for some kinds of faults and not others, every claim built on top of it inherits an unmeasured assumption.

The nearest existing work approaches this gap from adjacent but distinct angles, and each stops short of closing it. A benchmark of naturally-failing multi-agent trajectories with human-annotated responsible-agent and decisive-step labels establishes that attribution can be scored by agreement with human judgment, but agreement with a human is not the same claim as agreement with a manufactured cause. An intervention-driven auto-debugging system treats attribution as a testable causal hypothesis --- segmenting a session, proposing a suspected culprit, editing it, and replaying --- and argues explicitly that the standard single-decisive-culprit assumption is structurally weak, yet stops short of building the standardized, cross-component ground-truth corpus its own argument implies is needed. A tracer trained via reinforcement learning already injects faults into successful traces and replays them to manufacture labels by construction, the same trick this paper leans on --- but its injection targets one perturbed action to generate training supervision for a single tracer, not a controlled, taxonomy-spanning substitution protocol that reports a measured attribution-accuracy gap across fault types. The nearest neighbor of all is a concurrent large-scale corpus, Who\&When Pro \citep{WhoAndWhenPro2026}, that manufactures failed trajectories exactly the way this paper does --- replaying a successful prefix and injecting a fault --- at a scale (12{,}326 trajectories across three modalities) this paper does not attempt; but its ``golden label'' is assigned by convention at the moment of injection, with no matched no-op control, no identity-replay tier, no bypass arm, and no measurement of whether the observed failure tracks the injected change rather than stochasticity or the edit pathway itself. That unverified step is precisely the object of study here: the certification gate built below is the instrument such corpora implicitly assume, and its per-class validity measurement is the result none of them reports. A further cluster sits one level removed: benchmark-construction efforts that validate reproducibility by replaying a recorded trajectory, or that grade runs by final container state under a coarse execution/coherence/verification taxonomy, give the field a shared evaluation surface but no causal locus fixed by injection; step-level credit-assignment methods for single-model reasoning localize a step within one model's own chain of thought, but neither operate on a multi-component pipeline nor validate against an injected, known cause. And the domain most associated with fault ground truth, program-repair research, supplies a contrast rather than a template: a database of real, isolated software bugs certifies its ground truth backward, from a deterministic triggering test a candidate fix either passes or fails. No such test exists for a semantic plan-step error, a broken tool call, or a polluted observation in a natural-language agent trajectory --- ground truth here cannot be certified backward from a failure; it can only be manufactured forward, by injecting a known fault into a run that would otherwise have succeeded and checking whether the outcome actually moves.

Turning that forward-manufacturing idea into a working protocol surfaces its own difficulties. The three pipeline components are architecturally entangled --- the plan conditions what the executor attempts, which conditions which tools get called and what comes back --- so a correlational cue about which component looks anomalous is not causal evidence about which one is responsible, and no amount of passive observation resolves this on its own. The intervention meant to break that entanglement can introduce a confound of its own: a substituted plan can perturb the executor's behavior even with the executor unchanged, and ordinary stochasticity can flip an outcome independent of any injected fault, so the instrument itself needs a way to detect when it is the one manufacturing false positives. Building the fault library that makes this measurable at all is harder than it sounds --- injecting a fault does not guarantee the run fails for that reason, and a template can blur across categories in ways only an empirical audit, not a design review, can catch. And because a failure's observable symptom can sit many steps downstream of its true cause, with cascading effects making an innocent later component look responsible, an attribution method's accuracy has to be measured with enough resolution to expose exactly this misattribution, not hidden inside one pooled number.

Our approach turns each obstacle into a specific, checkable design choice. A checkpoint-and-replay harness holds two of the three components fixed at their known-good state while swapping in a faulted version of the third, breaking the entanglement by intervention rather than better post-hoc reasoning. A ground-truth label is certified only when the outcome actually flips against a matched no-op-substitution control run through the identical procedure --- so the substitution mechanism's own false-positive rate is measured, not assumed away. The fault library and its labels are audited against independent human judgment and contrasted against a naturally-failing, human-labeled comparison set, so realism is a reported quantity rather than a claimed property. Every candidate attribution method is then applied uniformly to the certified corpus and scored with a class-conditional confusion matrix against ground truth --- including a budget-matched control that gives a post-hoc judge the same replay budget an intervention-based method consumes, so an apparent advantage from intervention access can be told apart from one that is really just extra compute.

We evaluate this design on tool-use tasks drawn from an existing agent benchmark, contrasting the constructed, certified fault corpus against a naturally-failing, human-labeled comparison set, and characterizing a battery of existing attribution approaches --- post-hoc LLM-judge explanation, ablation-based blame, intervention-based do-then-verify, and a repurposed step-level localizer --- against the resulting ground truth, alongside trivial floor and compute-parity controls.

In summary, we contribute:
\begin{itemize}[leftmargin=1.4em]
\item A per-class certification criterion that turns injected faults into certified ground truth --- requiring an observed outcome flip against a matched no-op control, an identity-replay tier, and a fault-bypass arm --- rather than assigning labels by injection convention; and, as its direct corollary, the first per-class measurement of the substitution premise (``does failure track the substituted component?'') that injection-built corpora and intervention-based debuggers currently assume without testing.
\item A substitution-and-replay harness that isolates a single pipeline component's causal contribution to a task failure while explicitly measuring, rather than assuming away, the false-positive rate the substitution procedure itself might introduce.
\item A uniform, budget-matched evaluation battery --- the secondary readout --- that scores existing failure-attribution approaches against certified ground truth with a full class-conditional confusion matrix, isolating genuine causal reasoning from raw replay-compute advantage.
\item A released, reusable certified fault-trajectory corpus --- versioned per scaffold and backbone, with certification explicitly bounded to the scaffold and backbone it was manufactured on, and with a per-class label-survival rate reported whenever a tier is re-certified on a new backbone --- positioned to let future attribution methods be checked against the same certified labels rather than each earning trust through its own bespoke human-agreement study.
\end{itemize}

\section{Problem Definition}

\noindent\textbf{Formal setup.} A tool-use agent pipeline is organized around three component slots $K = \{plan, tool, obs\}$ together with a fixed executor $\varepsilon$. A task $t$, run under a fixed agent scaffold $A$, produces a run $R(t) \to (\tau, y)$: a trajectory $\tau$ and an outcome $y \in \{succ, fail\}$ determined by the task's own outcome test. A fault operator $\delta_c$ injects a known fault into component $c \in K$; $\delta_\emptyset^c$ denotes the class-matched no-op substitution on component $c$ --- a swap of $c$ for a semantically equivalent good copy via the \emph{same edit pathway} that class's fault operator uses, changing nothing about what the component does but reproducing that class's substitution disturbance exactly; $\delta_{id}$ denotes an identity replay with no edit at all, separating intrinsic replay stochasticity from edit-induced disturbance; and $\delta_{byp}^c$ denotes a bypass replay in which component $c$'s fault instrumentation is installed but its fault transformation is disabled --- the arm that ties an observed flip to the fault's content rather than to the instrumentation's presence. Given a base run that succeeds reproducibly, $k$ independent replays are executed from a checkpoint immediately after a substitution; $flip(t, \delta_c) = 1$ exactly when the base run is majority-success and the $k$ post-substitution replays are majority-failure. A ground-truth failure-locus label $\ell(\tau) = c$ is assigned only when $\delta_c$ was the applied substitution, $flip = 1$, and the bypass replay $\delta_{byp}^c$ restores majority-success, and only for fault classes $c$ whose flip rate under $\delta_c$ clears a pre-specified margin over the flip rate under its class-matched control $\delta_\emptyset^c$ --- classes that do not clear this margin are reported as uncertified rather than folded into a single global verdict. An attribution method is a function $m: \tau \to \hat{c}$, predicting one or more components as responsible; its correctness is judged against $\ell(\tau)$, restricted to the certified classes.

\noindent\textbf{Why this is hard.} The three components are causally entangled by construction --- the plan conditions what the executor attempts, which conditions which tools are invoked and what observations return --- so any cue correlated with a component's apparent involvement in a failure is not, on its own, evidence that the component is the true cause; only an intervention that holds two components fixed while perturbing the third can break this entanglement, and no amount of additional passively-observed trajectories resolves it. The intervention itself is not free of risk: substituting one component while nominally holding the others fixed can still perturb their behavior indirectly, and ordinary stochasticity in agent dialogue can flip $y$ independent of any injected fault --- so $flip(t, \delta_c)$ must be interpreted relative to $flip(t, \delta_\emptyset^c)$ and to the identity-replay tier $\delta_{id}$, not read as evidence of causation on its own; the two control tiers additionally disambiguate \emph{why} a control flips, since an edit-pathway artifact ($\delta_\emptyset^c$ flips, $\delta_{id}$ does not) is a harness defect to fix, not a finding; and the bypass arm $\delta_{byp}^c$ plays the complementary per-trial role, checking that removing the fault's content --- with the instrumentation still installed --- removes the failure. Constructing the fault operators $\{\delta_c\}$ that make $\ell(\tau)$ meaningful at scale is harder than a design exercise: injecting a fault does not guarantee $flip = 1$ for the intended reason, and a single fault template can blur across component categories, so unambiguity and realism are properties that must be measured empirically --- via the flip-rate margin and an independent audit --- rather than asserted from the generator's design. Finally, because a trajectory can be arbitrarily long, the observable symptom $y = fail$ may be causally and temporally distant from the injected component, and cascading effects can make an untouched downstream component appear responsible in $\tau$ even when $\ell(\tau)$ says otherwise, so an attribution method's accuracy must be measured conditionally on this distance rather than reported as one pooled number.

\section{Method Overview}

At a conceptual level, the study is a single measurement pipeline with five cooperating parts. A base task pool and a library of parameterized fault generators feed a substitution-and-replay harness; an outcome-flip gate turns the harness's raw replay outcomes into per-class-certified ground-truth labels; and a uniform, budget-matched application of every candidate attribution method against those certified labels, scored by a class-conditional confusion matrix, turns the whole exercise into a measured accuracy gap rather than an assumed one.

The \textbf{fault injection library} parameterizes three fault categories --- plan-step corruption, tool breakage, and observation pollution --- each emitting a fault template tied to a specific injected component and step index. The \textbf{substitution harness} checkpoints a reproducibly-succeeding run, holds two of the three components fixed at their known-good state, swaps in the faulted version of the third, and replays $k$ times from the injection point. The \textbf{outcome-flip certification gate} compares each fault class's replay outcomes against a matched control lattice --- class-matched no-op substitution, identity replay, and fault-bypass arm --- run through the identical procedure, and certifies a ground-truth label only for classes whose flip-rate gap clears a pre-specified margin and only for trials whose bypass replay restores success. The \textbf{uniform attribution battery} applies every candidate method --- post-hoc judges, ablation-based blame, intervention-based do-then-verify, and a repurposed step-level localizer --- to every certified trajectory under matched compute accounting, including a budget-matched judge variant that isolates causal-intervention access from raw replay compute. The \textbf{verification scorer} compares each method's predicted component against the certified ground truth with a full class-conditional confusion matrix, exposing not just an aggregate accuracy number but where each method's errors concentrate.

\section{Methodology}

\subsection{Fault Injection Library}
\textbf{Challenge.} No prior corpus scores a continuous, cross-component failure-locus label for tool-use agents at all, so injecting a fault does not by itself guarantee that unambiguity or realism hold --- a fault template can blur across component categories, and injecting a fault does not guarantee the run actually fails for the intended reason.

\textbf{Insight.} Realism and label-unambiguity cannot be guaranteed analytically by careful generator design; they can only be turned into a checkable empirical audit against independent human judgment, rather than an assumed property of the fault templates themselves.

\textbf{Method.} A parameterized generator library spans three fault categories --- plan-step deletion, reordering, or replacement; tool replacement with an error-returning, empty, or malformed-output stub; and observation-level noise injection, contradiction, or truncation --- each emitting a fault template bound to one injected component and one step index in a base trajectory. A sample of certified trajectories is independently audited by human raters, who are shown the trajectory without the injected-component label and asked to identify the responsible component themselves.

\subsection{Substitution Harness}
\textbf{Challenge.} The pipeline's three components are architecturally entangled, so a component that merely looks anomalous downstream of a failure is not necessarily the component actually responsible for it --- a confound that holds even under unlimited passive observation, because it is a property of the pipeline's structure rather than of the data's scale.

\textbf{Insight.} Entanglement cannot be resolved by more sophisticated post-hoc reasoning over the same observed trajectory; it can only be broken by an intervention that perturbs exactly one component while the checkpointed state of the others is replayed unchanged.

\textbf{Method.} A reproducibly-succeeding base run is checkpointed at a chosen injection point. Two of the three components ($\{plan, executor/tool, observation\}$ minus the target) are held fixed at their recorded good state; the faulted version of the third is swapped in via an orchestrator-message-level edit layer, preserving architectural generality across agent scaffolds; the run replays forward from the checkpoint for $k$ independent trials.

\subsection{Outcome-Flip Certification Gate}
\textbf{Challenge.} The substitution procedure itself is a candidate source of false positives --- a swapped component can perturb downstream behavior for reasons unrelated to the injected fault, and ordinary agent stochasticity can flip an outcome on its own, so an observed flip cannot automatically be trusted as evidence of the injected fault's causal effect.

\textbf{Insight.} Anchoring a ground-truth label to an \emph{observed effect}, relative to a matched control that undergoes the identical replay procedure with no injected fault, is what prevents the measurement instrument from silently mislabeling its own noise as a causal finding.

\textbf{Method.} For each fault class, the flip rate under the injected substitution is compared against the flip rate under a class-matched no-op substitution --- an equivalent good copy swapped in via the same edit pathway as that class's fault injection --- run through the identical checkpoint-and-replay procedure, alongside an identity-replay tier with no edit at all that separates edit-pathway disturbance from intrinsic replay stochasticity. A fourth, bypass arm re-runs the identical replay with the fault instrumentation installed but its transformation disabled; a certified trajectory must also recover majority-success under bypass, tying the flip to the injected fault's content rather than to the presence of the instrumentation. A fault class is certified only if the flip-rate gap over its no-op control clears a pre-specified margin; trajectories within a certified class are labeled with the injected component only when their own individual replay outcome flips and their bypass arm recovers. Uncertified classes are reported as such, not silently merged into the certified set.

\subsection{Uniform Attribution Battery}
\textbf{Challenge.} Any comparison across attribution methods that differ in how much information or intervention access they consume risks attributing an accuracy difference to causal reasoning when it is really explained by unequal resources.

\textbf{Insight.} Holding the trajectories, the certified labels, and the compute budget identical across methods is what makes a measured accuracy difference interpretable as a difference in reasoning quality rather than a difference in resources --- and adding a variant of the cheapest method that is deliberately given the expensive method's resource budget turns that interpretability concern into a direct, checkable comparison.

\textbf{Method.} Every candidate attribution method --- a post-hoc judge under multiple context-presentation strategies, an ablation/leave-one-out procedure that resets each component to its known-good state and checks which restoration flips the outcome back to success, an intervention-based do-then-verify procedure, and a step-level localizer repurposed from single-agent reasoning --- is applied to every certified trajectory. A budget-matched judge variant is additionally run: the same post-hoc judge, given the same number of additional sampled resamples of the failing run that the intervention-based method consumes as replay budget, but without the ability to actually intervene.

\subsection{Verification Scorer}
\textbf{Challenge.} A single pooled accuracy number can hide exactly the failure mode that matters most for a long, cascading trajectory: a method that performs well on faults adjacent to their symptom but poorly on temporally distant ones would look uniformly competent if only one aggregate number is reported.

\textbf{Insight.} Reporting the full class-conditional confusion structure, rather than one pooled number, is what turns ``how far below the ceiling are we'' and ``is the shortfall uniform across fault classes'' into two distinct, separately measurable readouts of the same scoring pass.

\textbf{Method.} Each method's predicted component is compared against the certified ground-truth label using exact-match top-1 accuracy, with a full three-by-three confusion matrix over $\{plan, tool, observation\}$, computed both overall and conditioned on fault class, with bootstrap confidence intervals reported throughout. Multi-component blame, where a method predicts more than one responsible component, is scored by precision and recall rather than exact match.

\section{Experiments}

\subsection{Experimental Setup}

\textbf{Datasets.} An existing GAIA-style tool-use task suite supplies the base task pool from which reproducibly-succeeding runs are selected for injection --- realistic multi-step tool-use scenarios are what makes the substitution protocol's findings relevant to deployed agents. The Who\&When dataset of naturally-failing, human-annotated multi-agent trajectories \citep{WhoAndWhen2025} serves as a contrast set: it lets this paper's construction-based labels be compared against human-annotated ones on the same kind of underlying failure, directly addressing the concern that injected faults might not represent the distribution of naturally-occurring failures. The constructed injected-fault corpus produced by this paper's own fault injection library and certification gate is the artifact every research question below is computed over. The corpus is released in versioned tiers keyed to scaffold and backbone: post-hoc methods are evaluable from the released trajectories alone, intervention-based methods require live checkpoint-replay and are therefore certified per backbone, and at least one full tier is built on a pinned open-weights backbone so re-certification and intervention baselines remain runnable after any API-hosted model is retired. Each re-certification of a tier on a new backbone reports a per-class \textbf{label-survival rate} --- the fraction of certified faults that still flip the outcome on the new backbone --- as a first-class corpus statistic, with non-survival treated as an expected, interpreted outcome (a stronger backbone can route around a weaker fault) rather than a defect of the corpus.

\textbf{Baselines} (sanity-check comparators, not competitors to beat --- this is a benchmark-construction and measurement study, and no research question here is framed as ``our procedure beats baseline X''). A post-hoc LLM-judge, applied under multiple context-presentation strategies \citep{WhoAndWhen2025}, is the standard explanation-based attribution approach being characterized. Ablation/leave-one-out blame and an intervention-based do-then-verify procedure \citep{DoVer2025} are the standard heuristic and interventionist comparators. A step-level localizer, repurposed from single-agent reasoning \citep{CRM2025,FAPO2025} to this multi-component setting, tests whether a method built for a different granularity transfers. A trained fault-attribution tracer \citep{AgenTracer2025} is included conditionally: only if its weights or prompts are publicly released; otherwise it is retained as the design precedent for this paper's own outcome-flip labeling trick and dropped from the accuracy table with a stated availability note. A random/majority-class control and a surface-heuristic floor --- blaming the component hosting the last step before the first observable error signal, where the signal is detected by a frozen, deterministic rule set over logged tool-return status fields, schema-invalid payloads, and fixed exception patterns, with no LLM call --- check whether the certified corpus carries any signal beyond a trivial guess. A budget-matched post-hoc judge isolates causal-intervention access from raw replay compute.

\textbf{Contamination control.} The base task pool is drawn from an existing, publicly released tool-use benchmark, so a post-hoc LLM-judge that has memorized a task's public solution path could look accurate for a reason unrelated to reasoning about the injected fault --- recognizing that a step deviates from a remembered correct trajectory is not the same as causally identifying which pipeline component the injected fault actually touched, and this alternative explanation is distinct from the compute-parity confound the budget-matched control already isolates. To rule it out, a held-out subset of at least 30 base tasks not present in the benchmark's public release (or task instances with paraphrased surface wording that breaks exact-string memorization) is scored separately from the main corpus; a method whose held-out accuracy falls 10 or more points below its full-corpus accuracy is flagged as relying on memorized task familiarity rather than genuine trajectory-level causal reasoning, and its headline accuracy is reported as memorization-assisted, with only the held-out accuracy treated as evidence of causal attribution ability.

\textbf{Piloting.} All arm sizes below are gated by a pre-registered walking-skeleton pilot --- a miniature of the full control lattice, on the order of twenty base tasks at reduced replay depth --- whose pass/fail thresholds on base-task yield, no-op flips, identity flips, and the tool-break flip-rate gap are fixed before any scale-up; RQ2's power simulation is run on this pilot, and a failed gate routes to a pre-named harness-fix or negative-result branch rather than ad hoc redesign.

\textbf{Metrics.} Human-audit agreement rate with the constructed ground-truth label, per fault category; per-class outcome-flip rate and its gap over the class-matched no-op control, plus the identity-replay flip rate and the bypass recovery rate; the per-class certification-rejection rate --- the fraction of injected, outcome-flipping trials that a convention labeler would accept but the control lattice rejects; top-1 component-attribution accuracy and the full plan/tool/observation confusion matrix, overall and per certified fault class, with base-task-cluster-bootstrap confidence intervals throughout --- trials sharing a base task form one cluster, so the effective sample size is the number of base tasks, not the number of trials; precision/recall when multi-component blame is scored; the held-out-subset accuracy gap used as the contamination check.

\textbf{Research questions:}
\begin{enumerate}[leftmargin=1.6em]
\item RQ1 --- Do independent human auditors agree with the constructed ground-truth failure-locus label?
\item RQ2 --- Does the observed task failure track the substituted pipeline component?
\item RQ3 --- Does the no-op control itself stay within a safe stochastic-flip bound?
\item RQ4 --- Do existing attribution methods beat a trivial surface-heuristic floor on the certified corpus?
\item RQ5 --- Does an intervention-based method's advantage survive giving a post-hoc judge an equal replay budget?
\item RQ6 --- How far below the by-construction ceiling does the best-performing method fall?
\item RQ7 --- Is the accuracy shortfall concentrated in plan-level faults specifically?
\end{enumerate}

\subsubsection{RQ1: Do independent human auditors agree with the constructed ground-truth failure-locus label?}
\textbf{Setup.} A sample of 200 certified trajectories, balanced across the three fault categories, is shown to independent human raters without the injected-component label; raters identify which component they judge responsible. Agreement with the constructed label is computed overall and per fault category.

\textbf{Expected outcome.} Human auditors agree with the constructed label for at least 80\% of the sampled trajectories, per fault category; agreement below this threshold in a given category triggers re-examination of that category's fault templates before any further release.

\textbf{Hypothesis (plain language).} The fault templates are unambiguous enough, in the specific sense that an independent human reading the trajectory would reach the same causal conclusion the construction procedure did.

\textbf{Metrics.} Human-label agreement rate, overall and per fault category.

\textbf{Result.} \ph

\subsubsection{RQ2: Does the observed task failure track the substituted pipeline component?}
\textbf{Setup.} For each of the three fault classes, at least 100 injection trials and at least 100 class-matched no-op-substitution trials (via the same edit pathway as that class's injection) are run, plus at least 100 identity-replay trials, each replayed $k = 4$ times from a checkpoint on a reproducibly-succeeding base run. A two-proportion comparison with cluster-robust inference at the base-task level (trials sharing a base task form one cluster) tests each injected class's flip rate against its class-matched control, one-sided, $\alpha = 0.05$, Benjamini-Hochberg corrected across the three classes; 95\% cluster-bootstrap confidence intervals are reported on all rates. The design is powered against a true gap of 0.30, not the 0.20 floor itself --- at the floor, the conjunctive pass criterion (point-estimate margin plus significance) succeeds only about half the time --- and a class whose simulated cluster-robust power on the pilot falls below 0.8 is pre-declared inconclusive rather than scored as null. In addition, the per-class certification-rejection rate is reported: the fraction of injected trials that a convention labeler --- one that accepts any injected, outcome-flipping trial, as injection-built corpora with convention labels do --- would admit but the control lattice (no-op margin, identity bound, bypass recovery) rejects. This statistic is computed from the same arms and costs no additional runs; it is the direct quantitative comparison between certified and convention-assigned labels.

\textbf{Expected outcome.} Each fault class's flip rate exceeds its class-matched no-op control's flip rate by at least 0.20 absolute, with the gap statistically significant after correction; a class that does not clear this margin is reported as uncertified rather than merged into the certified set, consistent with the possibility that some fault types (plan-level corruption in particular) propagate too diffusely to certify at this replay depth --- though a reversed pattern is reported at full prominence rather than smoothed into the expected wording.

\textbf{Hypothesis (plain language).} Failure tracking the substituted component is a real but potentially class-dependent phenomenon, not a premise that automatically holds uniformly across every kind of fault.

\textbf{Metrics.} Per-class outcome-flip rate and its gap over the class-matched no-op control, with 95\% cluster-bootstrap CIs; the per-class certification-rejection rate relative to a convention labeler.

\textbf{Result.} \ph

\subsubsection{RQ3: Does the no-op control itself stay within a safe stochastic-flip bound?}
\textbf{Setup.} The same class-matched no-op-substitution trials used in RQ2's control arms are examined on their own, alongside the identity-replay arm: their flip rates under $k = 4$ replays at temperature 0 are compared against pre-registered upper bounds --- 0.10 per class for the class-matched no-op arms and 0.05 for the identity-replay tier.

\textbf{Expected outcome.} Each class-matched no-op arm's flip rate remains at or below 0.10, and the identity-replay tier at or below 0.05. The two tiers disambiguate a violation: if the identity tier passes but a class-matched arm fails, the edit pathway itself is the artifact --- a harness defect to be fixed and re-run, never a finding; if the identity tier itself fails, the pipeline is intrinsically chaotically sensitive, undermining every certified label built on top of it, and the design allows one pre-registered replay-depth increase before pivoting to the global negative-result fallback.

\textbf{Hypothesis (plain language).} Ordinary agent stochasticity, on its own, does not manufacture enough false-positive outcome flips to corrupt the certification gate.

\textbf{Metrics.} Class-matched no-op flip rate per class and identity-replay flip rate, each with 95\% cluster-bootstrap CI.

\textbf{Result.} \ph

\subsubsection{RQ4: Do existing attribution methods beat a trivial surface-heuristic floor on the certified corpus?}
\textbf{Setup.} Computed only over the certified tier established by RQ2's per-class flip-rate check; an uncertified fault class is excluded from this pool rather than silently included. Each non-trivial method (post-hoc judge, ablation/leave-one-out, intervention-based do-then-verify, step-level localizer) is scored against the surface-heuristic floor via a McNemar-style test on paired correct/incorrect outcomes with base-task-cluster-robust inference (a cluster bootstrap of the paired accuracy difference, resampling base tasks), $\alpha = 0.05$, Benjamini-Hochberg corrected across the method-by-fault-category grid. If the surface heuristic itself already matches ground truth at 0.90 or above on the certified tier, this comparison is re-scored under a pre-registered heuristic-ceiling branch: the heuristic-sufficiency rate is reported per fault class as a finding in its own right --- the measured boundary of where attribution is trivial --- and method comparisons are restricted to the residual slice where the heuristic fails.

\textbf{Expected outcome.} Each non-trivial method family exceeds the floor's accuracy by at least 0.10 absolute on the certified corpus.

\textbf{Hypothesis (plain language).} The certified corpus carries a real, learnable attribution signal beyond what a trivial heuristic already captures --- otherwise the whole exercise of building and comparing methods against it would be measuring noise.

\textbf{Metrics.} Per-method accuracy and its gap over the floor, with 95\% cluster-bootstrap CIs.

\textbf{Result.} \ph

\subsubsection{RQ5: Does an intervention-based method's advantage survive giving a post-hoc judge an equal replay budget?}
\textbf{Setup.} Computed only over the certified tier. The budget-matched post-hoc judge variant --- given the same number of additional sampled resamples the intervention-based method consumes, without intervention access --- is compared against the standard intervention-based method via a base-task-cluster-robust paired bootstrap difference test on accuracy, Benjamini-Hochberg corrected alongside RQ4's grid.

\textbf{Expected outcome.} The accuracy gap between the intervention-based method and the budget-matched judge narrows to below 0.10 absolute, indicating that at least part of any raw advantage for intervention-based methods reflects extra compute rather than causal-intervention access specifically.

\textbf{Hypothesis (plain language).} An apparent advantage for methods that can intervene is not automatically evidence of better causal reasoning --- it needs to survive a fair compute comparison first.

\textbf{Metrics.} Accuracy gap between the intervention-based method and its budget-matched judge counterpart, with bootstrap CI.

\textbf{Result.} \ph

\subsubsection{RQ6: How far below the by-construction ceiling does the best-performing method fall?}
\textbf{Setup.} Computed only over the certified tier established by RQ2. The best-performing method's overall accuracy, across the certified fault classes, is estimated with a 95\% cluster-bootstrap confidence interval (10,000 resamples, resampling base tasks rather than trials). The surface-heuristic floor is excluded from the ``best method'' maximization; its accuracy is reported separately as the heuristic-sufficiency rate.

\textbf{Expected outcome.} The best-performing method's accuracy does not exceed 0.90, with the CI's upper bound below 0.95 --- a measurable gap remains below the ceiling of 1.0 that a by-construction ground truth makes possible in principle.

\textbf{Hypothesis (plain language).} Mechanical verification against a causally certified ground truth exposes a real accuracy shortfall in current attribution methods, rather than confirming they already perform at the achievable ceiling.

\textbf{Metrics.} Best-method overall accuracy with 95\% cluster-bootstrap CI; the heuristic-sufficiency rate reported separately.

\textbf{Result.} \ph

\subsubsection{RQ7: Is the accuracy shortfall concentrated in plan-level faults specifically?}
\textbf{Setup.} Computed only over the certified tier. For the post-hoc judge family, accuracy conditioned on plan-level faults is compared against accuracy conditioned on tool-level faults via a class-conditional cluster-bootstrap difference test (resampling base tasks), Benjamini-Hochberg corrected alongside RQ4's grid. Because certification keeps only outcome-flipping faults and could otherwise leave different surviving severity distributions per class, the comparison is run on severity-matched template pairs --- plan and tool templates at the same pre-registered severity level --- with the surface-heuristic floor's per-class accuracy reported alongside as a difficulty proxy.

\textbf{Expected outcome.} The post-hoc judge family's accuracy on plan-level faults is at least 0.10 lower than on tool-level faults, consistent with plan-corruption effects propagating more diffusely through the trajectory before manifesting as a visible failure. If the gap vanishes under severity matching, the claim is downgraded to a severity-composition effect rather than a fault-class effect; a significant reversed ordering is reported at full prominence rather than folded into a null wording.

\textbf{Hypothesis (plain language).} If current methods have a specific blind spot, it is more likely to be attributing blame correctly when the fault is a distant, semantic plan error than when it is a nearby, mechanically obvious tool break.

\textbf{Metrics.} Class-conditional accuracy gap (plan vs.\ tool) for the post-hoc judge family, with bootstrap CI.

\textbf{Result.} \ph

\vspace{0.5em}
\noindent\textbf{Further evaluation.} With more time, a natural extension is compound and multi-component faults --- trajectories where more than one component is simultaneously faulted, or where responsibility is genuinely distributed rather than single-locus --- explicitly out of scope here because the current certification gate and confusion-matrix scorer are built around a single injected component per trial, and would need their own labeling invariant rather than a direct extension of the outcome-flip gate above.

\section{Related Work}

\textbf{Failure attribution validated against human or single-culprit labels.} A benchmark of naturally-failing multi-agent trajectories \citep{WhoAndWhen2025} formalizes attribution as identifying decisive errors and releases a dataset of failure logs with human annotations of the responsible agent and step, framing attribution as an LLM-as-a-Judge task under several context-presentation strategies. An intervention-driven auto-debugging system \citep{DoVer2025} instead treats attribution as a testable causal hypothesis, segmenting a session, proposing a suspected culprit, and replaying from an edited checkpoint to check whether the outcome changes --- and argues explicitly that the standard single-decisive-culprit assumption is structurally weak. Both establish real, usable machinery for attribution --- the first a human-labeled comparison corpus this paper draws on directly, the second a rewind-edit-replay mechanism this paper's own substitution harness borrows --- but neither validates its own blame assignments against a controlled, causally manufactured ground truth spanning a plan/tool/observation taxonomy, which is the gap this paper's certification gate is built to close.

\textbf{By-construction fault injection as training supervision.} A trained attribution tracer \citep{AgenTracer2025} generates its own high-fidelity supervision by injecting faults into successful traces and replaying to check whether the outcome flips --- the closest real precedent for this paper's labeling trick, and the source of the inject-then-replay idea this paper's certification gate builds on. Its injection, however, targets a single perturbed action to produce step-level supervision for training a dedicated tracer, not a controlled, taxonomy-spanning substitution protocol that reports a measured, class-conditional attribution-accuracy gap across independently-evaluated methods --- the extension this paper makes.

\textbf{Injection-built corpora with convention labels.} The nearest neighbor to this paper is Who\&When Pro \citep{WhoAndWhenPro2026}, a concurrent large-scale benchmark that constructs 12{,}326 failed trajectories by replaying a successful prefix and injecting an erroneous action, labeling the injected step a ``golden label.'' It occupies the manufacture-forward territory at a scale this paper does not attempt, and this paper concedes that territory outright: injection-built ground truth for agent failure attribution is, by itself, no longer new. What Who\&When Pro does not do is certify. Its label is assigned by convention at the moment of injection --- no class-matched no-op control, no identity-replay tier, no bypass arm, no per-class flip-rate margin, and no measurement of whether the observed failure tracks the injected change rather than stochasticity, an edit-pathway artifact, or a dormant weakness of the base run. Its golden labels are, exactly, the hypothesis this paper's RQ2 and RQ3 test, and the certification-rejection rate reported there quantifies how many convention-accepted labels survive controlled scrutiny. The two artifacts are therefore complements --- scale without certification, and certification at moderate scale --- but the validity measurement, and the certified tier it licenses, are claimed here.

\textbf{Benchmark reproducibility and coarse failure taxonomies.} Agent-benchmark construction work has moved toward validating trajectories by replaying them to confirm reproducibility before accepting a task \citep{SelfEvolveBench2025}, and toward grading agent runs by final-state outcome tests annotated with a coarse execution/coherence/verification failure taxonomy \citep{TerminalBench2025}. Related work inducing interpretable behavioral metrics from open-ended human feedback \citep{AutoLibra2025} argues, independently, that coarse outcome-based metrics miss the intermediate decisions that actually cause success or failure. All three give the field a shared, higher-quality evaluation surface, but none fixes a ground-truth causal locus by injection --- a reproducible or well-taxonomized failure is not the same claim as a causally certified one. (A concurrent line of 2026 work applies a related trajectory-replay check to validate \emph{newly synthesized, harder} benchmark tasks rather than to attribute a \emph{pre-existing} failure to one of several candidate causal components \citep{SelfEvolveBench2025} --- the replay mechanism is a genuine precedent for this paper's own reproducibility discipline, but it answers a different question: is this evolved task well-posed, not which fixed pipeline component caused this failure.)

\textbf{Step-level credit assignment for single-agent reasoning, and the classical fault-localization contrast.} A survival-model reward formulation \citep{CRM2025} assigns each reasoning step a credit reflecting its causal contribution to the final outcome, and a separate reward-modeling approach \citep{FAPO2025} trains a detector for ``flawed-positive'' rollouts that also predicts the earliest erroneous step. Both localize a step within a single model's own chain of reasoning, but neither operates on a multi-component agent pipeline where the true fault may be a broken tool interface or a polluted observation rather than a bad token-level step, and neither validates its localization against an injected, known cause. The domain most associated with fault ground truth, classical software fault-localization research, supplies a useful contrast rather than a template: a database of real, isolated Java bugs \citep{Defects4J2014} certifies its ground truth backward, from a deterministic triggering test each candidate fix either passes or fails. No such test exists for a semantic plan-step error, a broken tool call, or a polluted observation in a natural-language agent trajectory --- which is precisely why this paper's ground truth is manufactured forward, by injection and observed outcome flip, rather than annotated backward from a naturally-occurring failure.

\textbf{Concurrent work.} Beyond Who\&When Pro, treated above as the nearest neighbor, a parallel line of 2026 work also pursues causal attribution for agent failures, but at a different granularity or along a different manipulated variable. Two efforts (``Causal Agent Replay,'' ``CausalFlow'') perform counterfactual intervention at the level of a single reasoning step within one trajectory --- a finer-grained target than this paper's fixed three-component taxonomy --- and report validation only against synthetic structural causal models with planted ground truth rather than real agent trajectories. A third (``POIROT''/BLAME) repurposes a multi-agent system's own constituent agents as a decentralized diagnostic layer, an architectural intervention on the evaluator rather than a controlled substitution of the pipeline component under test. None manipulates the same plan/tool/observation taxonomy under the same substitution-and-replay protocol this paper uses, so none is treated as a direct competitor here.

\section{Conclusion}

This work reframes a question every failure-attribution method for tool-use agents implicitly answers, and even the newest injection-built corpora assume rather than test: does blame assigned to a pipeline component --- or a label assigned at the moment of injection --- correspond to a cause, or only to a plausible-looking symptom? Concurrent work has already begun manufacturing ground truth forward at scale; what remains untested, and what this paper is built to measure, is whether the manufactured label deserves its name. The central methodological wager is that this question is only answerable through a certification lattice --- an injected fault whose outcome flip is required to clear a matched no-op control, an identity-replay bound, and a bypass-recovery check --- rather than through backward annotation, which is not an option available to natural-language agent pipelines the way it is to code-level software bugs, or through injection convention, which assumes the very premise under test. Whichever fault classes turn out to support certified ground truth, and however far current attribution methods fall short of the resulting ceiling, the finding is decision-relevant: a class that fails certification means substitution-based ground truth for that fault type needs a different protocol before it can be trusted, and a method that fails to beat the trivial floor means its previously-reported accuracy against human-agreement labels does not survive mechanical verification. Three limitations bound this claim from the outset and are natural next steps rather than afterthoughts: the design is scoped to single-component faults with a computable ground truth, leaving compound and distributed-responsibility failures for follow-up under their own labeling invariant; the certification gate's reliance on a reproducibility gate and bounded replay budget means results are reported per fault class rather than as one global verdict, which this paper treats as an honest constraint of the method rather than a gap to paper over; and certification is scaffold- and backbone-specific --- the released corpus certifies attribution on the scaffold and backbone it was manufactured on, with versioned per-backbone tiers and one pinned open-weights tier standing in for the cross-scaffold generality that a single release cannot claim and that only a future multi-scaffold replication could establish.

\vspace{1em}
\hrule
\vspace{1em}

\section*{Pre-handoff checklist}
\begin{itemize}[leftmargin=1.4em]
\item \textbf{Organization/navigation} --- narrative flows problem $\to$ method $\to$ experiments $\to$ related work without forward references to undefined terms.
\item \textbf{Terminology discipline} --- ``certified ground truth,'' ``outcome flip,'' ``substitution harness,'' ``control lattice,'' ``bypass arm,'' ``label-survival rate,'' ``certification-rejection rate,'' ``budget-matched control,'' ``confusion matrix'' each defined at first use in Problem Definition or Method Overview, used consistently thereafter; no internal stage IDs (\texttt{sp\_*}, \texttt{m\_*}, \texttt{rq\_*}, C1--C4, M1--M5) leaked into prose --- checked directly against the discussion pipeline log.
\item \textbf{Notation with worked example} --- the formal setup section defines $(K, \varepsilon, t, A, \tau, y, \delta_c, \delta_\emptyset^c, \delta_{id}, \delta_{byp}^c, k, flip, \ell, m)$ and the Problem Definition/Methodology sections consistently reuse this notation; no separate worked numeric example is included since no numbers exist yet (would require fabricating results).
\item \textbf{RQ framing matches outcomes} --- all 7 RQs are graded/quantitative (rates, accuracy gaps, agreement proportions), each answered in its own subsection with an explicit falsifiable expected outcome, not a binary yes/no.
\item \textbf{No typesetting artifacts} --- only explicit placeholder markers in Results, per the drafting convention this is a pre-registered spec, not a report.
\item \textbf{Artifact plan} --- the substitution harness, fault injection library, and certified fault-trajectory corpus are the natural release artifacts; not yet built, so not claimed as available.
\end{itemize}

\section*{Hand-off}
This is a draft, not a finished paper --- experiments have not run. Natural next steps: prose polish and venue-specific LaTeX formatting once a target venue is chosen, or revising an earlier stage if piloting surfaces a problem (e.g., if a fault class fails RQ2's certification margin entirely).

\bibliographystyle{plainnat}
\bibliography{refs}

\end{document}
